% =============================================================================
% GKB BEATRIX × MABE 2030 Proposal
% FehrAdvice & Partners AG
% Version 2.0 | Februar 2026
% =============================================================================

\documentclass[11pt,a4paper]{article}

% --- Packages ---
\usepackage[utf8]{inputenc}
\usepackage[T1]{fontenc}
\usepackage[ngerman]{babel}
\usepackage[margin=2.5cm, top=3cm, bottom=3cm]{geometry}
\usepackage{xcolor}
\usepackage{tabularx}
\usepackage{booktabs}
\usepackage{tcolorbox}
\usepackage{enumitem}
\usepackage{fancyhdr}
\usepackage{graphicx}
\usepackage{hyperref}
\usepackage{titlesec}
\usepackage{array}
\usepackage{colortbl}
\usepackage{lastpage}
\usepackage{ragged2e}

% --- FehrAdvice Farben ---
\definecolor{fadarkblue}{RGB}{2,64,121}
\definecolor{falightblue}{RGB}{84,158,222}
\definecolor{fadarkgray}{RGB}{37,33,42}
\definecolor{falightgray}{RGB}{243,245,247}
\definecolor{falilac}{RGB}{161,160,198}
\definecolor{famint}{RGB}{126,189,172}
\definecolor{faocher}{RGB}{222,203,63}
\definecolor{faorange}{RGB}{222,157,62}

% --- Typography ---
\usepackage{helvet}
\renewcommand{\familydefault}{\sfdefault}

% --- Section Formatting ---
\titleformat{\section}
  {\Large\bfseries\color{fadarkblue}}
  {\thesection.}{0.5em}{}
\titleformat{\subsection}
  {\large\bfseries\color{fadarkblue}}
  {\thesubsection}{0.5em}{}
\titleformat{\subsubsection}
  {\normalsize\bfseries\color{fadarkgray}}
  {\thesubsubsection}{0.5em}{}

% --- Header/Footer ---
\pagestyle{fancy}
\fancyhf{}
\fancyhead[L]{\small\color{fadarkgray}FehrAdvice \& Partners AG}
\fancyhead[R]{\small\color{fadarkgray}BEATRIX $\times$ MABE 2030}
\fancyfoot[L]{\small\color{fadarkgray}Vertraulich}
\fancyfoot[C]{\small\color{fadarkgray}Version 2.0 | Februar 2026}
\fancyfoot[R]{\small\color{fadarkgray}Seite \thepage\ von \pageref{LastPage}}
\renewcommand{\headrulewidth}{0.4pt}
\renewcommand{\footrulewidth}{0.4pt}

% --- Hyperlinks ---
\hypersetup{
  colorlinks=true,
  linkcolor=fadarkblue,
  urlcolor=falightblue,
  citecolor=fadarkblue,
}

% --- Custom tcolorbox styles ---
\tcbuselibrary{skins,breakable}

\newtcolorbox{fabox}[1][]{
  colback=falightgray,
  colframe=fadarkblue,
  fonttitle=\bfseries\color{white},
  coltitle=white,
  colbacktitle=fadarkblue,
  #1
}

\newtcolorbox{fahighlight}[1][]{
  colback=fadarkblue!5!white,
  colframe=fadarkblue,
  boxrule=0.5pt,
  #1
}

\newtcolorbox{faaccent}[1][]{
  colback=famint!10!white,
  colframe=famint!80!black,
  boxrule=0.5pt,
  #1
}

% --- Table row colors ---
\newcommand{\headerrow}{\rowcolor{fadarkblue}}
\newcommand{\altrow}{\rowcolor{falightgray}}

% =============================================================================
\begin{document}

% --- Title Page ---
\thispagestyle{empty}
\begin{center}
\vspace*{3cm}

{\color{fadarkblue}\rule{\textwidth}{2pt}}

\vspace{1.5cm}

{\Huge\bfseries\color{fadarkblue} BEATRIX $\times$ MABE 2030}

\vspace{0.8cm}

{\Large\color{fadarkgray} Verhaltensökonomische Marktbearbeitungsstrategie\\[0.3cm] mit KI-Unterstützung}

\vspace{2cm}

{\color{fadarkblue}\rule{0.4\textwidth}{0.5pt}}

\vspace{1.5cm}

\begin{tabular}{rl}
\textbf{\color{fadarkblue}Für:} & Graubündner Kantonalbank (GKB), Geschäftsleitung \\[0.3cm]
\textbf{\color{fadarkblue}Von:} & FehrAdvice \& Partners AG \\[0.3cm]
\textbf{\color{fadarkblue}Datum:} & Februar 2026 \\[0.3cm]
\textbf{\color{fadarkblue}Version:} & 2.0 \\[0.3cm]
\textbf{\color{fadarkblue}Bezug:} & MABE Schlüsselereignisse Briefing
\end{tabular}

\vfill

{\small\color{fadarkgray} FehrAdvice \& Partners AG | Klausstrasse 20 | 8008 Zürich | fehradvice.com}

\vspace{0.5cm}
{\color{fadarkblue}\rule{\textwidth}{2pt}}

\end{center}
\newpage

% =============================================================================
% 1. AUFTRAGSVERSTÄNDNIS
% =============================================================================

\section{Auftragsverständnis}

Die GKB will ihre Marktbearbeitungsstrategie 2030 auf \textbf{Schlüsselereignisse der Generation X} ausrichten --- differenziert nach Kohorten, fundiert durch Daten, umsetzbar im Graubündner Kontext (Briefing, Punkte a--f).

Wir lesen das als zwei aufeinander aufbauende Fragen:

\begin{center}
\begin{tabularx}{\textwidth}{c >{\raggedright\arraybackslash}p{4.5cm} X}
\toprule
\headerrow \textcolor{white}{\textbf{Etappe}} & \textcolor{white}{\textbf{Kernfrage}} & \textcolor{white}{\textbf{Ergebnis}} \\
\midrule
1 & \textbf{Was} sind die richtigen Schlüsselereignisse --- und \textbf{für wen}? & Priorisierte Schlüsselereignisse, differenziert nach Teilzielgruppen und Motiv-Typen \\
\altrow 2 & \textbf{Wie} bespielt die GKB diese Schlüsselereignisse? & Kommunikation, Produktportfolio-Mapping, Kampagnenplanung \\
\bottomrule
\end{tabularx}
\end{center}

\vspace{0.3cm}

Unsere Rolle: Die verhaltensökonomische Fundierung beider Fragen. Schlüsselereignisse sind nicht nur demografische Datenpunkte --- sie sind \textbf{Entscheidungssituationen}, in denen Kund:innen handeln oder eben nicht handeln. Wir helfen zu verstehen, was in diesen Momenten passiert und wie die GKB sie besser nutzen kann.

% =============================================================================
% 2. WIE WIR ARBEITEN
% =============================================================================

\section{Wie wir arbeiten: das Fragen-Abo}

\begin{fabox}[title=Das Prinzip]
Jeden Monat stellt die GKB eine strategische Frage --- oder zwei operative. Wir beantworten sie vollständig: mit Daten, Modell, Simulation und GL-ready Report. So entsteht Monat für Monat ein vollständiges Bild. Reihenfolge, Tempo und Priorisierung bestimmt die GKB.
\end{fabox}

\vspace{0.3cm}

\begin{center}
\begin{tabularx}{\textwidth}{l X}
\toprule
\headerrow \textcolor{white}{\textbf{Kontingent}} & \textcolor{white}{\textbf{Pro Monat}} \\
\midrule
\textbf{1 strategische Fragestellung} & z.B. \guillemotleft Welche Schlüsselereignisse sind für GenX prioritär?\guillemotright \\
\altrow \textbf{oder 2 operative Fragestellungen} & z.B. \guillemotleft Wie reagiert Kohorte 1 auf Eigenmietwert-Reform?\guillemotright\ + \guillemotleft Kampagne Segment 50+?\guillemotright \\
\bottomrule
\end{tabularx}
\end{center}

\subsection{Wie wir eine Frage bearbeiten}

\begin{center}
\begin{tabularx}{\textwidth}{>{\bfseries\color{fadarkblue}}l X}
\toprule
\headerrow \textcolor{white}{\textbf{Schritt}} & \textcolor{white}{\textbf{Was passiert}} \\
\midrule
GKB liefert & Daten und Kontext zur Fragestellung (CRM, Hypotheken, Berater-Einschätzungen) \\
\altrow Wir kalibrieren & Daten-Onboarding, GR-spezifische Kontextfaktoren laden \\
BEATRIX modelliert & Verhaltensökonomische Analyse, Simulation, Parametrisierung \\
\altrow Wir liefern & GL-ready Report mit konkreten Empfehlungen, 1 Präsentation \\
\bottomrule
\end{tabularx}
\end{center}

\subsection{Investition}

\begin{center}
\begin{tabularx}{\textwidth}{X r}
\toprule
\headerrow \textcolor{white}{\textbf{Komponente}} & \textcolor{white}{\textbf{Investition}} \\
\midrule
\textbf{Monatliches Abo} (1 strategische oder 2 operative Fragen) & CHF 12'000/Mt. \\
\altrow \textbf{Mindestlaufzeit} & 3 Monate \\
\textbf{Startinvestition (3 Monate)} & \textbf{CHF 36'000} \\
\bottomrule
\end{tabularx}
\end{center}

\vspace{0.2cm}

\textbf{Inklusive:} Alle Deliverables als GL-ready Reports, 1 Präsentation/Monat (vor Ort oder remote), Reisekosten Chur, Zugang zum GKB-Kontextvektor (350 Faktoren, laufend aktualisiert). \textbf{Monatlich kündbar} nach den ersten 3 Monaten.

\subsection{Team}

\begin{center}
\begin{tabularx}{\textwidth}{l X}
\toprule
\headerrow \textcolor{white}{\textbf{Rolle}} & \textcolor{white}{\textbf{Verantwortung}} \\
\midrule
\textbf{Senior Partner} & Gesamtsteuerung, GL-Präsentationen \\
\altrow \textbf{Behavioral Economist} & Schlüsselereignis-Parametrisierung, Modellierung \\
\textbf{Data Analyst} & GKB-Datenanalyse, Share-of-Wallet, Segmentierung \\
\altrow \textbf{Prof.\ Ernst Fehr} & Wissenschaftliche Begleitung, Quality Gate \\
\bottomrule
\end{tabularx}
\end{center}

% =============================================================================
% 3. FRAGESTELLUNGEN
% =============================================================================

\section{Fragestellungen auf Basis eures Briefings}

Aus dem Briefing (Punkte a--f) leiten wir folgende Fragestellungen ab. Im Kick-off priorisieren wir gemeinsam.

\begin{center}
\begin{tabularx}{\textwidth}{c X c >{\raggedright\arraybackslash}p{2cm}}
\toprule
\headerrow \textcolor{white}{\textbf{\#}} & \textcolor{white}{\textbf{Fragestellung}} & \textcolor{white}{\textbf{Bezug}} & \textcolor{white}{\textbf{Typ}} \\
\midrule
F1 & Welche Schlüsselereignisse und Bedürfnisräume sind relevant --- priorisiert nach Potenzial und GKB-Stärken? & a, b & Strategisch \\
\altrow F2 & Welche Teilzielgruppen und Motivationstypen gibt es --- jenseits der Alterskohorte? & c, d & Strategisch \\
F3 & Welche Verhaltenshürden verhindern, dass Kund:innen bei diesen Ereignissen handeln? & b, c & Strategisch \\
\altrow F4 & Welche Massnahmen wirken --- und in welcher Reihenfolge? & e & Strategisch \\
F5 & Wie positioniert sich die GKB gegenüber Wettbewerbern pro Schlüsselereignis? & f & Strategisch \\
\altrow F6 & GL-Entscheidungsvorlage: Prioritäten, Investment, ROI & Gesamt & Operativ \\
F7 & Kommunikation für eine Zielkohorte oder einen Motivationstyp & e & Operativ \\
\altrow F8 & MABE-2030-Fahrplan & Gesamt & Operativ \\
\bottomrule
\end{tabularx}
\end{center}

\subsection{Beispiel-Fahrplan (erste 3 Monate)}

\begin{center}
\begin{tabularx}{\textwidth}{c >{\raggedright\arraybackslash}p{3cm} X >{\raggedright\arraybackslash}p{3.5cm}}
\toprule
\headerrow \textcolor{white}{\textbf{Monat}} & \textcolor{white}{\textbf{Typ}} & \textcolor{white}{\textbf{Beispiel-Fragestellung}} & \textcolor{white}{\textbf{Deliverable}} \\
\midrule
1 & 1 strategische & F1: Schlüsselereignisse priorisieren & Priorisiertes Ranking mit Business Case \\
\altrow 2 & 1 strategische & F2: Teilzielgruppen \& Motivationstypen & Kohorten- und Typenprofile \\
3 & 1 strat.\ oder 2 oper. & F3: Verhaltenshürden --- oder 2 operative Fragen & Barriers Map oder 2 oper.\ Deliverables \\
\bottomrule
\end{tabularx}
\end{center}

\vspace{0.1cm}
\textit{Reihenfolge, Tempo und Priorisierung bestimmt die GKB.}

\vspace{0.3cm}

\begin{faaccent}
\textbf{Nach 3 Monaten} hat die GKB eine solide Entscheidungsgrundlage. Das Abo kann danach monatlich verlängert werden, z.B.\ für Interventions-Design, Pilotbegleitung, neue Kohorten oder Geschäftskundenanalyse.
\end{faaccent}

% =============================================================================
% 4. ERSTE ERKENNTNISSE
% =============================================================================

\section{Erste Erkenntnisse}

Die folgenden Beobachtungen basieren auf dem Briefing, öffentlich zugänglichen Daten und unserer bisherigen GKB-Vorarbeit. Sie zeigen, was wir bereits sehen --- und wo die vertiefte Analyse mit GKB-Daten ansetzt.

\subsection{Die 3 Kohorten sind grundverschieden}

\begin{center}
\begin{tabularx}{\textwidth}{>{\raggedright\arraybackslash}p{2.5cm} X X X}
\toprule
\headerrow & \textcolor{white}{\textbf{Vor-Pensionäre}} & \textcolor{white}{\textbf{Pensionsplaner}} & \textcolor{white}{\textbf{Vermögensaufbauer}} \\
\midrule
\textbf{Jahrgänge} & 1965--1969 (56--60 J.) & 1970--1974 (51--55 J.) & 1975--1980 (45--50 J.) \\
\altrow \textbf{Leitfrage} & \guillemotleft Wie sichere ich meinen Übergang?\guillemotright & \guillemotleft Wie schliesse ich die Lücke?\guillemotright & \guillemotleft Wie baue ich systematisch auf?\guillemotright \\
\textbf{Modus} & Absicherung + Optimierung & Planung + Aufstockung & Aufbau + Automatisierung \\
\altrow \textbf{Zentraler Bias} & Time Inconsistency & Ostrich Effect & Present Bias \\
\bottomrule
\end{tabularx}
\end{center}

\vspace{0.2cm}

\textbf{Kohorte 1} steht vor irreversiblen Entscheidungen (Rente vs.\ Kapital, Frühpensionierung). \textbf{Kohorte 2} weiss, dass sie handeln müsste, vermeidet aber die Konfrontation mit der Rentenlücke. \textbf{Kohorte 3} braucht Automatisierung --- nicht \guillemotleft machen Sie etwas\guillemotright, sondern \guillemotleft wir haben es für Sie eingerichtet\guillemotright.

\subsection{Graubünden ist kein Durchschnitts-Markt}

\begin{center}
\begin{tabularx}{\textwidth}{>{\raggedright\arraybackslash}p{3.5cm} >{\raggedright\arraybackslash}p{4cm} X}
\toprule
\headerrow \textcolor{white}{\textbf{Besonderheit}} & \textcolor{white}{\textbf{Fakten}} & \textcolor{white}{\textbf{Implikation}} \\
\midrule
Überalterung & 65+ Anteil GR 23\% vs.\ CH 19\% & Pensionierungsberatung früher und häufiger relevant \\
\altrow Zweitwohnungen ($\sim$50\%) & $\sim$80'000, oft mit Hausbank ZKB/UBS & GKB als \guillemotleft GR-Expertin\guillemotright\ positionieren \\
Eigenmietwert-Reform & GR-Verlust $\sim$CHF 90 Mio/J. & 2025--2028 = natürliches Beratungsfenster \\
\altrow Tourismus-KMU & Saisonale Liquiditätsmuster & Relevant für Geschäftskunden und BVG \\
\bottomrule
\end{tabularx}
\end{center}

\vspace{0.2cm}

Eine 1:1-Übernahme nationaler Strategien funktioniert nicht. Die MABE muss GR-spezifisch sein --- und genau das ist der Kern von Frage F1.

\subsection{Was bereits vorliegt: GKB-Kontextvektor}

FehrAdvice hat im Rahmen der bestehenden Kundenbeziehung bereits erhebliche Vorarbeit geleistet:

\begin{center}
\begin{tabularx}{\textwidth}{l c X}
\toprule
\headerrow \textcolor{white}{\textbf{Layer}} & \textcolor{white}{\textbf{Faktoren}} & \textcolor{white}{\textbf{Beispiele}} \\
\midrule
\textbf{FIN} Financial & 70 & Geschäftsvolumen CHF 76.1 Mrd, Kommissionsertrag +7.7\% \\
\altrow \textbf{MKT} Market & 50 & 120'000 Privatkund:innen, 45 Standorte \\
\textbf{STR} Strategy & 35 & Hybrid-Strategie, Diversifikation Kommissionsgeschäft \\
\altrow \textbf{PEO} People & 60 & $\sim$1'000 MA, 12 J.\ Betriebszugehörigkeit \\
\textbf{TEC} Technology & 45 & GKB Mobile, Gioia 3a/Kids, Hypomat \\
\altrow \textbf{ORG} Governance & 40 & Kantonales Gesetz, FINMA-Compliance \\
\textbf{RIS} Risk & 30 & CET-1 18.8\%, S\&P AA, Staatsgarantie \\
\altrow \textbf{STK} Stakeholder & 20 & Kanton GR 84.6\%, SIX-kotiert \\
\bottomrule
\end{tabularx}
\end{center}

\vspace{0.2cm}

4 Use Cases sind bereits vormodelliert, mit kalibrierten Verhaltensparametern ($\lambda_\text{GKB} = 2.0$, $\tau_\text{TRUST} = 0.85$, $\beta_\text{REGIONAL} = 0.90$). BEATRIX kalibriert und verfeinert das bestehende Modell mit GKB-internen Daten.

\vspace{0.3cm}

\begin{fahighlight}
\textbf{Hinweis:} Alle Einschätzungen in diesem Kapitel sind ein erster Wurf. Die Priorisierung (Briefing: \guillemotleft aufgrund Potenzial und Stärken der GKB\guillemotright) erfordert GKB-interne Daten und ist Kern der ersten Monate.
\end{fahighlight}

% =============================================================================
% 5. NÄCHSTE SCHRITTE
% =============================================================================

\section{Nächste Schritte}

\begin{center}
\begin{tabularx}{\textwidth}{X l l}
\toprule
\headerrow \textcolor{white}{\textbf{Schritt}} & \textcolor{white}{\textbf{Timing}} & \textcolor{white}{\textbf{Verantwortung}} \\
\midrule
Proposal besprechen & Donnerstag & FehrAdvice + GKB \\
\altrow GKB-interne Abstimmung & 1--2 Wochen & GKB \\
Auftragserteilung + Kick-off & Nach GL-Entscheid & Beidseitig \\
\altrow Datenvorbereitung (CRM-Export) & Vor Kick-off & GKB IT \\
Projektstart & T + 2 Wochen & FehrAdvice \\
\bottomrule
\end{tabularx}
\end{center}

\vspace{0.5cm}

\subsection{Warum FehrAdvice + BEATRIX}

\begin{fahighlight}
\begin{description}[leftmargin=0.5cm, style=nextline]
\item[\textcolor{fadarkblue}{Wissenschaft + Technologie}] Prof.\ Ernst Fehr, einer der weltweit meistzitierten Ökonomen. BEATRIX macht 50 Jahre Forschung (2'347 Papers, 153 Theorien, 852 Cases) operativ nutzbar.

\item[\textcolor{fadarkblue}{Graubünden-Expertise}] 350 GKB-Kontextfaktoren bereits modelliert. 4 Use Cases voranalysiert. Das Modell kennt Graubünden.

\item[\textcolor{fadarkblue}{Messbare Ergebnisse}] Jede Empfehlung kommt mit Wahrscheinlichkeit und Konfidenzintervall. Nach 3 Monaten hat die GL alles für eine datenbasierte Entscheidung.
\end{description}
\end{fahighlight}

% =============================================================================
% ANHANG: BEATRIX
% =============================================================================

\newpage
\section*{Anhang: BEATRIX auf einen Blick}
\addcontentsline{toc}{section}{Anhang: BEATRIX auf einen Blick}

\begin{fabox}[title=Was ist BEATRIX?]
\textbf{B}ehavioral \textbf{E}conomics \textbf{A}I \textbf{T}echnology for \textbf{R}esearch and \textbf{I}mplementation e\textbf{X}cellence

\vspace{0.3cm}

BEATRIX ist kein Chatbot. BEATRIX ist ein \textbf{Entscheidungsunterstützungssystem}, das menschliches Verhalten im Kontext berechnet.
\end{fabox}

\vspace{0.5cm}

\subsection*{Die 3 Schichten}

\begin{center}
\begin{tabularx}{\textwidth}{l l X}
\toprule
\headerrow \textcolor{white}{\textbf{Schicht}} & \textcolor{white}{\textbf{Funktion}} & \textcolor{white}{\textbf{Beispiel}} \\
\midrule
1.\ Formale Berechnung & Deterministische EBF-Gleichungen & $\theta_B = \theta_A \times \prod_i M(\Delta\Psi_i)$ \\
\altrow 2.\ Parameter-Store & 64+ validierte Verhaltensparameter & $\lambda = 2.0$, $\beta = 0.90$, $\gamma = -0.68$ \\
3.\ KI-Übersetzung & Formale Ergebnisse in natürliche Sprache & \guillemotleft Die Verlustaversion ist\ldots\guillemotright \\
\bottomrule
\end{tabularx}
\end{center}

\vspace{0.5cm}

\subsection*{Der zentrale Unterschied}

\begin{center}
\begin{tcolorbox}[colback=fadarkblue!5!white, colframe=fadarkblue, width=0.85\textwidth]
\begin{center}
\textbf{Traditionelle KI} generiert Text aus Wahrscheinlichkeiten.\\[0.2cm]
\textbf{BEATRIX} berechnet Verhalten aus Kontextfaktoren.\\[0.4cm]
\textit{Text ist beliebig reproduzierbar.\\
Verhaltensvorhersagen sind falsifizierbar.}
\end{center}
\end{tcolorbox}
\end{center}

\vspace{0.5cm}

\subsection*{Wissenschaftliche Basis}

\begin{center}
\begin{tabular}{lr}
\toprule
\textbf{Ressource} & \textbf{Umfang} \\
\midrule
Kuratierte Papers & 2'347 \\
Formalisierte Theorien & 153 \\
Dokumentierte Cases & 852 \\
Verhaltensökonomische Parameter & 64+ \\
Kontextfaktoren (Schweiz + GKB) & 754 \\
\bottomrule
\end{tabular}
\end{center}

\vfill

\begin{center}
{\color{fadarkblue}\rule{0.5\textwidth}{0.5pt}}\\[0.5cm]
{\small\color{fadarkgray} FehrAdvice \& Partners AG | Klausstrasse 20 | 8008 Zürich | fehradvice.com}
\end{center}

\end{document}
